\documentclass[aspectratio=169,12pt,t]{beamer}
\usepackage{graphicx}
\setbeameroption{hide notes}
\setbeamertemplate{note page}[plain]
\usepackage{listings}
\usepackage{eepic}

% header.tex: boring LaTeX/Beamer details + macros

% get rid of junk
\usetheme{default}
\beamertemplatenavigationsymbolsempty
\hypersetup{pdfpagemode=UseNone} % don't show bookmarks on initial view
\usepackage[export]{adjustbox}


% font
\usepackage{fontspec}
\setsansfont
  [ ExternalLocation = fonts/ ,
    UprightFont = *-regular ,
    BoldFont = *-bold ,
    ItalicFont = *-italic ,
    BoldItalicFont = *-bolditalic ]{texgyreheros}
\setbeamerfont{note page}{family*=pplx,size=\footnotesize} % Palatino for notes
% "TeX Gyre Heros can be used as a replacement for Helvetica"
% I've placed them in fonts/; alternatively you can install them
% permanently on your system as follows:
%     Download http://www.gust.org.pl/projects/e-foundry/tex-gyre/heros/qhv2.004otf.zip
%     In Unix, unzip it into ~/.fonts
%     In Mac, unzip it, double-click the .otf files, and install using "FontBook"

% named colors
\definecolor{offwhite}{RGB}{255,250,240}
\definecolor{gray}{RGB}{155,155,155}
\definecolor{purple}{RGB}{177,13,201}
\definecolor{green}{RGB}{46,204,64}

\definecolor{background}{RGB}{255,255,255}
\definecolor{foreground}{RGB}{24,24,24}
\definecolor{title}{RGB}{27,94,134}
\definecolor{subtitle}{RGB}{22,175,124}
\definecolor{hilit}{RGB}{197,75,140}
\definecolor{princolorA}{RGB}{122,0,128}
\definecolor{princolorB}{RGB}{0,106,207}
\definecolor{princolorC}{RGB}{26,154,34}
\definecolor{princolorD}{RGB}{225,103,27}
\definecolor{princolorE}{RGB}{17,17,17}
\definecolor{vhilit}{RGB}{255,0,128}
\definecolor{codehilit}{RGB}{255,0,128}
\definecolor{lolit}{RGB}{95,95,95}
\definecolor{myyellow}{rgb}{1,1,0.7}
\definecolor{nhilit}{RGB}{128,0,128}  % hilit color in notes
\definecolor{nvhilit}{RGB}{255,0,128} % vhilit for notes

\newcommand{\princolorA}{\color{princolorA}}
\newcommand{\princolorB}{\color{princolorB}}
\newcommand{\princolorC}{\color{princolorC}}
\newcommand{\princolorD}{\color{princolorD}}
\newcommand{\princolorE}{\color{princolorE}}
\newcommand{\hilit}{\color{hilit}}
\newcommand{\vhilit}{\color{vhilit}}
\newcommand{\nhilit}{\color{nhilit}}
\newcommand{\nvhilit}{\color{nvhilit}}
\newcommand{\lolit}{\color{lolit}}

% use those colors
\setbeamercolor{titlelike}{fg=title}
\setbeamercolor{subtitle}{fg=subtitle}
\setbeamercolor{institute}{fg=lolit}
\setbeamercolor{normal text}{fg=foreground,bg=background}
\setbeamercolor{item}{fg=foreground} % color of bullets
\setbeamercolor{subitem}{fg=lolit}
\setbeamercolor{itemize/enumerate subbody}{fg=lolit}
\setbeamertemplate{itemize subitem}{{\textendash}}
\setbeamerfont{itemize/enumerate subbody}{size=\footnotesize}
\setbeamerfont{itemize/enumerate subitem}{size=\footnotesize}

% page number
\setbeamertemplate{footline}{%
    \raisebox{5pt}{\makebox[\paperwidth]{\hfill\makebox[20pt]{\lolit
          \scriptsize\insertframenumber}}}\hspace*{5pt}}

% add a bit of space at the top of the notes page
\addtobeamertemplate{note page}{\setlength{\parskip}{12pt}}

% default link color
\hypersetup{colorlinks, urlcolor={hilit}}

\lstset{language=bash,
        basicstyle=\ttfamily\scriptsize,
        frame=single,
        commentstyle=,
        backgroundcolor=\color{offwhite},
        showspaces=false,
        showstringspaces=false
        }


% a few macros
\newcommand{\bi}{\begin{itemize}}
\newcommand{\bbi}{\vspace{24pt} \begin{itemize} \itemsep8pt}
\newcommand{\ei}{\end{itemize}}
\newcommand{\be}{\begin{enumerate}}
\newcommand{\bbe}{\vspace{24pt} \begin{enumerate} \itemsep8pt}
\newcommand{\ee}{\end{enumerate}}
\newcommand{\sbi}{\begin{itemize} \fontsize{9pt}{9.5}\selectfont}
\newcommand{\sbe}{\begin{enumerate} \fontsize{9pt}{9.5}\selectfont}
\newcommand{\ig}{\includegraphics}
\newcommand{\subt}[1]{{\footnotesize \color{subtitle} {#1}}}
\newcommand{\ttsm}{\tt \small}
\newcommand{\ttfn}{\tt \footnotesize}
\newcommand{\figh}[2]{\centerline{\includegraphics[height=#2\textheight]{#1}}}
\newcommand{\figw}[2]{\centerline{\includegraphics[width=#2\textwidth]{#1}}}
\newcommand{\fighboxed}[2]{\centerline{\includegraphics[frame,height=#2\textheight]{#1}}}
\newcommand{\figwboxed}[2]{\centerline{\includegraphics[frame,width=#2\textwidth]{#1}}}
\newcommand{\lfigh}[2]{\includegraphics[height=#2\textheight]{#1}}
\newcommand{\lfigw}[2]{\includegraphics[width=#2\textwidth]{#1}}


%%%%%%%%%%%%%%%%%%%%%%%%%%%%%%%%%%%%%%%%%%%%%%%%%%%%%%%%%%%%%%%%%%%%%%
% end of header
%%%%%%%%%%%%%%%%%%%%%%%%%%%%%%%%%%%%%%%%%%%%%%%%%%%%%%%%%%%%%%%%%%%%%%

% title info
\title{generic HMM for multi-parent populations}
\author{\vspace*{-10pt} \href{https://kbroman.org}{Karl Broman}}
\institute{Biostatistics \& Medical Informatics, UW{\textendash}Madison}
\date{\href{https://twitter.com/kwbroman}{\tt \scriptsize \color{foreground} @kwbroman}
\\[-4pt]
\href{https://kbroman.org}{\tt \scriptsize \color{foreground} kbroman.org}
\\[-4pt]
\href{https://github.com/kbroman}{\tt \scriptsize \color{foreground} github.com/kbroman}
\\[-4pt]
{\scriptsize \href{https://kbroman.org/Talk_GenericHMM}{\tt kbroman.org/Talk\_GenericHMM}}
}


\begin{document}

% title slide
{
\setbeamertemplate{footline}{} % no page number here
\frame{
  \titlepage

  \vfill \hfill
  \href{https://creativecommons.org/publicdomain/zero/1.0/}{\includegraphics[height=7mm]{Figs/cc-zero.png}}

  \note{
    These are slides for a talk for the CTC ({\tt
    www.complextrait.org/ctc2021/}) on 1 Sept 2021.

    Slides:
    \href{https://kbroman.org/Talk_GenericHMM/generic_hmm.pdf}{\tt
      kbroman.org/Talk\_GenericHMM/generic\_hmm.pdf} \\[12pt]
    Slides with notes:
    \href{https://kbroman.org/Talk_GenericHMM/generic_hmm_notes.pdf}{\tt
      kbroman.org/Talk\_GenericHMM/generic\_hmm\_notes.pdf} \\[12pt]
    Source: \href{https://github.com/kbroman/Talk_GenericHMM}{\tt github.com/kbroman/Talk\_GenericHMM}
    \\[24pt]
    Related paper on bioRxiv: \href{https://doi.org/gswx}{\tt doi.org/gswx}
  }

} }




\begin{frame}[c]{}

  \begin{columns}

    \column{0.5\textwidth}

    \centering
    {\color{title} Recombinant Inbred Lines}

    \bigskip

    \figh{Figs/ri8.pdf}{0.8}


    \column{0.5\textwidth}

    \centering
    {\color{title} Advanced Intercross Population}

    \bigskip

    \figh{Figs/hs.pdf}{0.8}

  \end{columns}


  \note{
    Multi-parent populations are experimental crosses starting from
    multiple inbred founder lines.

    Major examples include the Collaborative Cross, a set of 8-way
    recombinant inbred lines in mouse, and Hetereogeneous Stock, which
    have been developed in both mice and rats and are advanced
    intercross populations derived from 8 founders. The Diversity
    Outbred mouse population is similar to HS. In plants, multi-parent
    recombinant inbred lines are called MAGIC lines (for multiparent
    advanced generation inter-cross).

    The offspring chromosomes will be mosaics of the founder
    chromosomes. Multi-parent populations can be homozygous (like RIL)
    or heterozygous (like HS). The number of founders need not be 8.
  }

\end{frame}




\begin{frame}[c]{Genome reconstruction}


  \figh{Figs/genome_reconstr.pdf}{0.9}

  \note{

    A key step in the analysis of multi-parent populations is genome
    reconstruction: using dense SNP genotypes in the founders and MPP
    offspring to infer the haplotypes across the genome.

    Here we consider a 1 Mbp region on chromosome 14 in a single
    Diversity Outbred Mouse. Open and closed circles indicate AA and
    BB genotypes at SNPs. Gray circles indicate AB heterozygous
    genotypes. Using the SNP data along the chromosome, we can
    calculate the probability of each possible genotype at each
    position.

    For this mouse, the left half of the interval looks to be
    homozygous DD, while the right half looks to be heterozygous AD.
  }

\end{frame}



\begin{frame}[c]{QTL genome scan}


  \figh{Figs/qtl_scan.png}{0.85}

  \note{

    One could skip the whole genome reconstruction and just do QTL
    analysis at the SNPs, as is done in GWAS.
    If the genotyped SNPs include individual causal polymorphisms,
    this could be best.

    But if there are multiple causal polymorphisms in a region QTL
    analysis with the inferred haplotypes may be more powerful.
    Moreover, if the founder strains have been sequenced, you can use
    the reconstructed genomes to get inferred genotypes at all
    polymorphisms in the founders. (Similar approaches were used in
    human GWAS, based on HapMap SNPs.)

    Here, the single-SNP analysis shows significant evidence for a
    single QTL on chromosome 1. The haplotype analysis indicates
    evidence for a second QTL on chromosome 4.

    Beyond QTL mapping, genome reconstructions are useful in data
    diagnostics. For example, the estimated number of crossovers is
    useful when assessing sample quality.
  }

\end{frame}



\begin{frame}[c]{DO genome}


  \figh{Figs/do_genome.pdf}{0.85}

  \note{
    Here is the reconstructed genome of a Diversity
    Outbred mouse. (The white segments are undetermined.)

    Our goal is to figure this out, using SNP genotypes on this mouse
    plus the 8 founder lines.
  }

\end{frame}





\begin{frame}{Hidden Markov model}

\figw{Figs/hmm.pdf}{1.0}

\bigskip

{
\centering
\renewcommand{\arraystretch}{2.0}

\begin{tabular}{l@{\hspace{1cm}}l}
Initial    & $\pi(g) = \text{Pr}(G_1 = g)$ \\
Transition & $t_i(g,g') = \text{Pr}(G_{i+1} = g' \ | \ G_i = g)$ \\
Emission   & $e_i(g) = \text{Pr}(O_i \ | \ G_i = g)$
\end{tabular}

}



  \note{
    The main approach for genome reconstruction is to use a hidden
    Markov model. The underlying diplotypes we're trying to determine
    follow a Markov chain $\{G_i\}$, but are unobserved. We observe
    SNP genotypes $\{O_i\}$, with an assumed conditional independence
    structure, where given $G_i$, $O_i$ is conditionally independent
    of everything else.

    Three sets of parameters govern the model: the initial and
    transition probabilities, which concern the pattern of underlying
    genotypes on the MPP chromosomes; and the emission probabilities,
    which relate the underlying genotypes to the observed SNP
    genotypes and largely concern a model for SNP genotyping errors.
  }

\end{frame}





\begin{frame}[c]{Exact probabilities}

  \only<1|handout 0>{\figw{Figs/broman2005.png}{0.8}}

  \only<2>{

    \begin{columns}

      \column{0.5\textwidth}

      \figw{Figs/broman2005.png}{0.95}

      \bigskip

      \figw{Figs/teuscher2007.png}{0.95}

      \column{0.5\textwidth}

      \figw{Figs/broman2012a.png}{0.95}

      \bigskip

      \figw{Figs/broman2012b.png}{0.95}

      \end{columns}
  }

  \note{
      I've spent quite a lot of time studying the pattern of genotypes
      on MPP chromosomes, first with a paper on multi-way recombinant
      inbred lines, but then following up with three further papers
      considering extra generations of outbreeding, the genotypes at
      intermediate generations, and the patterns in advanced
      intercross populations such as Diversity Outbred mice.

      The mathematics is interesting but tedious. And is it necessary?
      It would be nice to have a generic approach that could be used
      generally.
  }

\end{frame}





\begin{frame}[c]{Generic model}

  \begin{columns}

    \column{0.5\textwidth}

    \figh{Figs/founder_pop.pdf}{0.5}


    \column{0.5\textwidth}

    $k$ founders in proportions $\{\alpha_i\}$ \\[4pt]
    $n$ generations of random mating \\[18pt]
    {\color{title} Random chromosome}: \\[4pt]
    \qquad $\pi_i = \alpha_i$ \\[4pt]
    \qquad $t_{ij} = \alpha_j \ [\text{1} - (\text{1} - r)^n] \quad
    \text{when } i \ne j$ \\[18pt]
    {\color{title} Map expansion:} \\[4pt]
    \qquad $n (\text{1}-\sum\alpha_i^{\text{2}})$ \\[4pt]
    \qquad $= n \left(\frac{k-\text{1}}{k}\right) \quad \text{if } \alpha_i \equiv \text{1}/k$


  \end{columns}


  \note{
    And that is what I propose here. Imagine a population of $k$
    founders in known (but not necessarily equal) proportions, and
    that a multi-parent population is formed by random mating for $n$
    discrete generations. In this case, we can calculate the
    transition probabilities exactly.

    We could apply these equations more generally. We need just specify the
    proportions of the founders (which should be known from the design
    of the study) and the effective number of generations of random
    mating. The latter might be calibrated by considering the map
    expansion (the proportional increase in the number of
    recombination breakpoints, relative to a single meiosis).
    This could be approximated by computer simulation.

    For a heterozygous population, like HS or the DO, we draw two
    random chromosomes. For a homozygyous population, like MAGIC lines
    or the Collaborative Cross, we can pretend that they are doubled
    haploids, with a single random chromosome like above.

    For the X chromosome, we can use the same equations, replacing $n$
    with $\left(\frac{\text{2}}{\text{3}}\right) n$, due to
    recombination only happening on the X chromosome in females.
  }

\end{frame}






\begin{frame}{DO application}

  \figh{Figs/do_qtl.pdf}{0.78}

  \bigskip

  \hfill {\footnotesize \lolit data from Al-Barghouthi et al (2021)
    \href{https://doi.org/gkf64n}{\tt \lolit doi.org/gkf64n}}

  \note{
    If we apply our approach to data from Diversity Outbred mice, the
    results with the generic model proposed above are basically
    identical to the use of the more-exact model.  For data from
    Al-Bargouthi et al (2021), this is the biggest difference seen:
    the LOD curves are not distinguishable, as the biggest difference
    is just 0.01.
  }

\end{frame}






\begin{frame}[c]{CC038 X chr}

  \figh{Figs/cc_xchr_reconstr.pdf}{0.9}

  \note{
    Reconstructions of the genomes of Collaborative Cross lines are
    identical for autosomes, but there are important differences for
    the X chromosome.

    This slide shows the reconstruction of the X chromosome in
    Collaborative Cross line CC038, but the exact model (top panel)
    and by the approximate model (bottom panel).

    The analysis differs in that the top model excludes three of the
    eight founders and weighs one of the other five more highly.

    These results differ in a region around 135~Mbp, where in the bottom
    panel, B6 and NOD are assigned equally probability, as they are
    identical in the region, but the top panel was able to exclude B6.
  }

\end{frame}





\begin{frame}[c]{X chr in CC}

  \figh{Figs/ri8X.pdf}{0.95}

  \note{
    The X chromosome in the Collaborative Cross behaves different than
    autosomes. We list the crosses female $\times$ male; note that the
    Y chromosome comes from the H strain and the X chromosome comes
    from the five strains A, B, C, E, and F, with the average
    proportion from the C strain between twice that of the others.

    This can be really useful information (provided that it is
    correct), particularly as the X chromosome shows reduced
    polymorphism compared to the autosomes. Many of the CC
    founders share large stretches of DNA on the X chromosome.
  }

\end{frame}





\begin{frame}[c]{Summary}

  \bbi

\item Generic model for genome reconstruction in multi-parent
  populations

\item Specify relative proportions of founders \\
  + effective number of generations of random mating

\item Basic conclusion: {\hilit HAPPY is effective}

\item Implemented in \href{https://kbroman.org/qtl2}{R/qtl2} as cross types {\tt genril{\emph n}} and
  {\tt genail{\emph n}} \\
  (replacing {\tt \emph n} with the number of founders)

\item bioRxiv manuscript:
  \href{https://doi.org/gswx}{\tt \lolit doi.org/gswx}


  \ei

  \note{
    It's always good to provide a summary.
  }

\end{frame}



\begin{frame}[c]{}

\Large

Slides: \href{https://kbroman.org/Talk_GenericHMM}{\tt kbroman.org/Talk\_GenericHMM}

\vspace*{-7mm}
\hfill
\href{https://creativecommons.org/publicdomain/zero/1.0/}{\includegraphics[height=7mm]{Figs/cc-zero.png}}

\vspace{3mm}

bioRxiv manuscript:
\href{https://doi.org/gswx}{\tt \lolit doi.org/gswx}

\vspace{4mm}

\href{https://kbroman.org}{\tt \lolit kbroman.org}

\vspace{4mm}

\href{https://github.com/kbroman}{\tt \lolit github.com/kbroman}

\vspace{4mm}

\href{https://twitter.com/kwbroman}{\tt \lolit @kwbroman}

\vspace{4mm}

\href{https://kbroman.org/qtl2}{\tt \lolit kbroman.org/qtl2}




\note{
  Here's where you can find me and these slides, as well as a preprint
  giving further details on the work.
}

\end{frame}






\end{document}
